\begin{circuitikz}[scale=0.8, every node/.style={transform shape}, font=\sffamily\small]
% ---------------- the adapter
\draw[rounded corners=2pt, fill=extcol] (0,2.4) rectangle (3.6,6.6);
\node[font=\sffamily\bfseries\small] at (1.8,6.2) {Renkforce USB/TTL};
\node[font=\sffamily\scriptsize] at (1.8,5.85) {3.3 V logic};
\node[draw=linecol, fill=white, minimum width=16mm, minimum height=4mm,
      font=\sffamily\scriptsize] at (1.8,2.7) {jumper = 3V3};
\draw[thick] (1.8,6.6) -- (1.8,7.4);
\node[anchor=south, font=\sffamily\scriptsize] at (1.8,7.45) {to a Pi 3 USB-A port};
\node[anchor=east, font=\sffamily\scriptsize] at (3.45,5.0) {3V3};
\node[anchor=east, font=\sffamily\scriptsize] at (3.45,4.4) {GND};
\node[anchor=east, font=\sffamily\scriptsize] at (3.45,3.8) {TXD};
\node[anchor=east, font=\sffamily\scriptsize] at (3.45,3.2) {RXD};
% ---------------- the target
\draw[rounded corners=2pt, fill=fwcol] (9.0,2.4) rectangle (13.2,6.6);
\node[font=\sffamily\bfseries\small] at (11.1,6.2) {ESP32 NodeMCU};
\node[font=\sffamily\scriptsize] at (11.1,5.85) {target, 3.3 V pins};
\node[anchor=west, font=\sffamily\scriptsize] at (9.15,5.0) {VIN / 3V3};
\node[anchor=west, font=\sffamily\scriptsize] at (9.15,4.4) {GND};
\node[anchor=west, font=\sffamily\scriptsize] at (9.15,3.8) {TX};
\node[anchor=west, font=\sffamily\scriptsize] at (9.15,3.2) {RX};
% ---------------- four wires, transmit crossed to receive
\draw (3.6,5.0) -- (9.0,5.0);
\draw (3.6,4.4) -- (9.0,4.4);
\draw (3.6,3.8) -- (5.4,3.8) -- (7.2,3.2) -- (9.0,3.2);
\draw (3.6,3.2) -- (5.4,3.2) -- (7.2,3.8) -- (9.0,3.8);
\node[anchor=south, font=\sffamily\scriptsize] at (6.3,5.05) {3.3 V from the adapter};
\node[anchor=south, font=\sffamily\scriptsize] at (6.3,4.45) {common ground, fitted first};
\node[anchor=north, font=\sffamily\scriptsize, align=center] at (6.3,2.95)
  {transmit crossed to receive};
% ---------------- ESP32 download mode
\draw (13.2,4.8) -- (14.4,4.8);
\draw (14.4,4.8) to[nos] (14.4,3.6) node[ground]{};
\node[anchor=west, font=\sffamily\scriptsize, align=left] at (14.6,4.8) {GPIO0 low\\= download};
\draw (13.2,2.8) -- (14.4,2.8);
\node[anchor=west, font=\sffamily\scriptsize] at (14.6,2.8) {EN: reset};
% ---------------- the other target and its own UART
\draw[rounded corners=2pt, fill=fwcol] (9.0,-1.6) rectangle (13.2,1.0);
\node[font=\sffamily\bfseries\small] at (11.1,0.6) {ESP8266-PROG};
\node[font=\sffamily\scriptsize] at (11.1,0.2) {second target};
\draw[dashed] (1.8,2.4) -- (1.8,-0.3) -- (9.0,-0.3);
\node[anchor=south, font=\sffamily\scriptsize] at (5.8,-0.25)
  {or the same four wires to this board instead, one target at a time};
\draw[thick] (11.1,-1.6) -- (11.1,-2.8);
\draw[red!70!black, line width=1pt] (10.85,-2.45) -- (11.35,-1.95);
\draw[red!70!black, line width=1pt] (10.85,-1.95) -- (11.35,-2.45);
\node[anchor=north, font=\sffamily\scriptsize, align=center] at (11.1,-2.9)
  {onboard USB-UART};
\node[anchor=east, font=\sffamily\scriptsize, text=red!70!black, align=right] at (10.6,-2.2)
  {not while the\\TTL cable is on};
\end{circuitikz}
